\chapter*{Conclusion générale et Perspectives}
\addcontentsline{toc}{chapter}{Conclusion générale et Perspectives}
\markboth{Conclusion générale et Perspectives}{}

Ce mémoire vise la détection des émotions multiples dans les textes,
une tâche où les approches mono-étiquette et les représentations
textuelles plates peinent à refléter la cooccurrence des affects et
le déséquilibre sévère des classes sur des corpus tels que
GoEmotions \cite{demszky2020}.

L'architecture \textbf{HAKE-MER} a été conçue de manière modulaire
afin de mesurer, à protocole constant sur GoEmotions, l'apport de
chaque brique. Les expériences menées sur DistilBERT-base (trois
graines) montrent que \textbf{M2} (attention croisée et têtes par
émotion) porte le gain principal en F1-macro test ($\approx 0{,}505$),
que \textbf{M1} seul sous-performe une baseline PLM flat
($\approx 0{,}486$), et que \textbf{M3} (NRC ou SenticNet) n'améliore
pas de façon robuste la macro par rapport à M1+M2 (H3 tranchée au
chapitre~\ref{chap:experimentation}). Il reste à évaluer \textbf{M4}
(H4) sur la pile M3 retenue, puis à compléter l'analyse par classe et
les exemples qualitatifs.

Les prolongements possibles (autres corpus, graphes de connaissances,
modèles généralistes) sont discutés au chapitre~\ref{chap:perspectives}.
